\begin{annexitem}{annex:p18-example-all-pass-bound}{thm:p18-example-all-pass-bound}{The independent-sampling bound and its limitation}
\begin{proof}
Let \(X_i\) be the success indicator of trial i. Independence and equal
Bernoulli success probabilities give
\(P(X_1=\cdots=X_n=1)=\prod_{i=1}^{n}P(X_i=1)=p^n\).
Under \(p\leq1-\varepsilon\), monotonicity of the nth power on
nonnegative numbers gives \(p^n\leq(1-\varepsilon)^n\).
The specified procedure rejects the condition only on that all-pass event,
so this is exactly its false-rejection probability bound under the condition.

It is not a deductive population conclusion: for any
\(0<p\leq1-\varepsilon\), the all-pass event still has positive
probability p-to-the-n. Thus an admissible model can realize that observation
while the population condition remains true. Nor has a distribution over p
been supplied; a posterior probability of a claim about p cannot be calculated
from this sampling distribution alone. If selection depends on already known
success, the stated Bernoulli/independence-of-outcome model has not been
established, so its product formula does not apply to that selection process.
\end{proof}
\end{annexitem}

\begin{annexitem}{annex:p18-sign-bayes-bridge}{thm:p18-sign-bayes-bridge}{The declared probabilistic sign calculation}
\begin{proof}
The events H and its complement partition the sample space. By conditional
probability, \(P(S\cap H)=s\pi\) and
\(P(S\cap\neg H)=f(1-\pi)\). Add the probabilities of those disjoint
events to obtain \(P(S)=s\pi+f(1-\pi)\), positive by hypothesis.
Dividing \(P(S\cap H)\) by this value proves the displayed posterior.
Similarly \(P(\neg H\mid S)=f(1-\pi)/P(S)\). If f is positive,
the latter is positive because \(1-\pi>0\), so division gives
posterior odds \((s/f)(\pi/(1-\pi))\).

Multiply the inequality for posterior at least tau by the positive
denominator, then collect its two s-pi terms. The result is exactly
\((1-\tau)s\pi\geq\tau f(1-\pi)\); every step is reversible.
When s=f>0, cancel the common factor in the posterior expression,
obtaining pi. These arguments include s=0 or f=0 whenever the stated
denominator condition holds; the odds statement alone explicitly requires
f>0. They establish only identities in the declared probability model.
They do not infer causal direction, a correct reference class, calibrated
likelihoods, or a normative dialogue verdict from those identities.
\end{proof}
\end{annexitem}

\begin{annexitem}{annex:p18-sign-confounding-countermodel}{thm:p18-sign-confounding-countermodel}{Counting the finite process and both interventions}
\begin{proof}
There are two hundred equally probable input triples: two choices of C,
ten of U-X and ten of U-Y. At high load there are ninety X-triples,
of which fifty-four fail, and ten non-X triples, of which six fail.
At low load there are ten X-triples, of which one fails, and ninety
non-X triples, of which nine fail. These counts follow by multiplying
the allowed independent input choices in the two defining equations.

Therefore there are one hundred X-triples with fifty-five failures,
and one hundred non-X triples with fifteen failures. Their conditional
failure probabilities are respectively \(0.55\) and \(0.15\). Within high
load both conditional rates are \(0.6\); within low load both are \(0.1\).

On either intervention the failure equation is unchanged. Exactly sixty
of the one hundred high-load triples and ten of the one hundred low-load
triples satisfy it, giving \(70/200=0.35\). The installation input and its
replaced equation do not enter that equation, so the same count applies
whether X is set present or absent. This establishes the two distinct
intervention probabilities rather than substituting a conditional
observational rate for either of them. Hence positive association alone
cannot entail a positive causal intervention effect.
\end{proof}
\end{annexitem}
